\documentclass[11pt]{article}
\newcommand{\DocNumber}{VOLUME 07}
\newcommand{\DocTitle}{C++23 Clause-by-Clause Study Guide}
\newcommand{\DocSubtitle}{Thirty-three guided clause reviews and worked exercises}
\newcommand{\DocSummary}{A complete clause-coverage study map with misconception checks, answers, and standards-trace exercises.}
% Shared, pdfLaTeX-only design system for the C++23 documentation suite.
\usepackage[T1]{fontenc}
\usepackage[utf8]{inputenc}
\usepackage{lmodern}
\usepackage[english]{babel}
\usepackage[letterpaper,margin=0.82in,headheight=16pt,headsep=17pt]{geometry}
\usepackage{microtype}
\usepackage{amsmath,amssymb}
\usepackage{array,booktabs,tabularx,longtable}
\usepackage[table]{xcolor}
\usepackage[most]{tcolorbox}
\usepackage{listings}
\usepackage{fancyhdr}
\usepackage{titlesec}
\usepackage{enumitem}
\usepackage{hyperref}
\usepackage{needspace}
\usepackage{etoolbox}
\definecolor{ink}{HTML}{13253B}
\definecolor{teal}{HTML}{006E7F}
\definecolor{muted}{HTML}{5A6C7C}
\definecolor{fog}{HTML}{EDF4F6}
\definecolor{sand}{HTML}{FAF3E8}
\definecolor{rule}{HTML}{C9D7DF}
\definecolor{codebg}{HTML}{F7F9FB}
\definecolor{codecomment}{HTML}{47715D}
\hypersetup{colorlinks=true,linkcolor=teal,urlcolor=teal,pdfauthor={Independent C++23 engineering documentation},pdftitle={C++23 Technical Suite}}
\setlength{\parindent}{0pt}
\setlength{\parskip}{6pt}
\setcounter{tocdepth}{2}
\renewcommand{\arraystretch}{1.18}
\titleformat{\section}{\Large\bfseries\color{ink}}{\thesection}{0.7em}{}
\titleformat{\subsection}{\large\bfseries\color{teal}}{\thesubsection}{0.55em}{}
\titlespacing*{\section}{0pt}{17pt}{7pt}
\titlespacing*{\subsection}{0pt}{11pt}{4pt}
\pagestyle{fancy}
\fancyhf{}
\fancyhead[L]{\footnotesize\bfseries\color{teal}C++23 / ISO 14882:2024}
\fancyhead[R]{\footnotesize\color{muted}\DocNumber}
\fancyfoot[L]{\scriptsize\color{muted}Independent engineering companion, not an ISO publication}
\fancyfoot[R]{\scriptsize\color{muted}\thepage}
\renewcommand{\headrulewidth}{0.35pt}
\lstdefinestyle{cppstyle}{language=C++,basicstyle=\ttfamily\footnotesize,keywordstyle=\color{teal}\bfseries,commentstyle=\color{codecomment},stringstyle=\color{ink},backgroundcolor=\color{codebg},frame=single,rulecolor=\color{rule},breaklines=true,breakatwhitespace=true,columns=fullflexible,keepspaces=true,showstringspaces=false,tabsize=2,numbers=none,upquote=true,xleftmargin=3pt,xrightmargin=3pt,aboveskip=6pt,belowskip=8pt}
\lstset{style=cppstyle}
\newtcolorbox{standardbox}{enhanced,breakable,colback=fog,colframe=teal!35!white,boxrule=0.55pt,arc=1.4mm,left=3mm,right=3mm,top=2mm,bottom=2mm,before skip=7pt,after skip=8pt}
\newtcolorbox{practicebox}{enhanced,breakable,colback=sand,colframe=ink!20!white,boxrule=0.55pt,arc=1.4mm,left=3mm,right=3mm,top=2mm,bottom=2mm,before skip=7pt,after skip=8pt}
\newcommand{\anchorref}[1]{\par\smallskip{\footnotesize\textbf{Standard locator:} \texttt{#1}}\par}
\newcommand{\standard}[1]{\begin{standardbox}\textbf{Standard-grounded interpretation.} #1\end{standardbox}}
\newcommand{\practice}[1]{\begin{practicebox}\textbf{Engineering recommendation (not a normative ISO rule).} #1\end{practicebox}}
\newcommand{\takeaway}[1]{\par\textbf{Review question.} #1\par}
\newcommand{\docfront}{%
\begin{titlepage}
\vspace*{25mm}
{\color{teal}\bfseries\large ENGINEERING REFERENCE SERIES\par}
\vspace{12mm}
{\raggedright\color{ink}\bfseries\fontsize{24}{30}\selectfont\DocTitle\par}
\vspace{7mm}
{\color{muted}\fontsize{15}{19}\selectfont\DocSubtitle\par}
\vspace{12mm}
\noindent\color{teal}\rule{\linewidth}{2pt}\par
\vspace{9mm}
{\large\DocSummary\par}
\vfill
\begin{tabularx}{\linewidth}{@{}lX@{}}
\textbf{Volume} & \DocNumber\\
\textbf{Standards basis} & ISO/IEC 14882:2024 (seventh edition, C++23)\\
\textbf{Document type} & Original interpretive and instructional material\\
\textbf{Date} & 16 October 2026\\
\textbf{Toolchain} & pdfLaTeX; illustrative C++23 examples require separate compiler support\\
\end{tabularx}
\vspace{8mm}
{\footnotesize This document is not an authorized copy of, replacement for, or amendment to ISO/IEC 14882:2024. The original standard controls in the event of disagreement.}\par
\end{titlepage}
\pagenumbering{roman}
\tableofcontents
\clearpage
\pagenumbering{arabic}
\section*{How to use this volume}\addcontentsline{toc}{section}{How to use this volume}
The reference is \emph{ISO/IEC 14882:2024, Programming languages --- C++}. Section locators and square-bracketed stable names identify where to verify precise wording; they are not quotations. \textbf{Standard-grounded interpretation} boxes summarize language or library semantics. \textbf{Engineering recommendation} boxes contain optional design, security, portability, or process guidance. Code is original and illustrative; individual library facilities may not yet be implemented by every vendor. This suite does not grant rights to reproduce or distribute the licensed standard.\par\medskip
}
\newcommand{\docend}{\clearpage\section*{Reference and traceability}\addcontentsline{toc}{section}{Reference and traceability}
\textbf{Primary source.} ISO/IEC 14882:2024(en), \emph{Programming languages --- C++}, seventh edition, October 2024. Use its numbered clauses, stable names, Annex C (informative compatibility), and Annex D (compatibility features) for authoritative details.\par
\textbf{Limits.} Clause-level summaries and exercises in this document are selective explanations, not exhaustive reproductions of all semantic constraints, library mandates, complexity clauses, and overload participation requirements. A successful compiler test demonstrates behavior of that specific implementation and configuration, not universal conformance.\par
\textbf{Use of contemporary engineering tools.} Compilers, sanitizers, analyzers, CI systems, repository policies, and vulnerability taxonomies are implementation guidance rather than language-standard requirements.\par}

\begin{document}
\docfront
\section{How to study each numbered clause}
Each card is a study prompt, not a copied clause. Read the cited clause, identify a normative guarantee or limitation, explain it in your own words, produce a positive example, and then test a boundary condition. Preserve the standard's square-bracketed stable names when recording notes; numeric section labels can change between editions.

\small
\begin{longtable}{@{}p{.18\linewidth}p{.40\linewidth}p{.34\linewidth}@{}}
\toprule
\textbf{Pass} & \textbf{Technique} & \textbf{Evidence}\\
\midrule
\endfirsthead
\toprule
\textbf{Pass} & \textbf{Technique} & \textbf{Evidence}\\
\midrule
\endhead
1. Vocabulary & Extract controlled terms & One precise definition and counterexample\\
\addlinespace[2pt]
2. Semantics & Translate normative text into obligations & Assumptions and conclusion\\
\addlinespace[2pt]
3. Examples & Construct valid/invalid cases & Compile/run matrix and expected outcome\\
\addlinespace[2pt]
4. Implementation & Review portability and extensions & Vendor/target qualification\\
\addlinespace[2pt]
5. Recall & Explain without the source open & Answer and supporting locator\\
\addlinespace[2pt]
\bottomrule
\end{longtable}
\normalsize
\section{Clause 1: Scope}
\anchorref{1}
The scope identifies the subject of the language standard, not a security or style policy.

\standard{This clause belongs to ISO/IEC 14882:2024; its precise constraints must be verified in the standard, including subclauses and any cross references.}
\textbf{Common trap.} Misreading a coding guide as part of the normative language standard.

\textbf{Self-test.} What kind of claim is outside the standard even if it is good engineering?

\textbf{Answer rationale.} An organization-specific code review gate.

\practice{Make a one-page note with the relevant stable names, one original well-formed example, a negative or edge case, and a recorded toolchain assumption.}
\section{Clause 2: Normative references}
\anchorref{2}
Referenced documents are needed in the appropriate context to apply the specification.

\standard{This clause belongs to ISO/IEC 14882:2024; its precise constraints must be verified in the standard, including subclauses and any cross references.}
\textbf{Common trap.} Assuming every non-ISO webpage has normative force.

\textbf{Self-test.} How should a conflicting external tutorial be handled?

\textbf{Answer rationale.} Consult the normative source and its defined dependency.

\practice{Make a one-page note with the relevant stable names, one original well-formed example, a negative or edge case, and a recorded toolchain assumption.}
\section{Clause 3: Terms and definitions}
\anchorref{3}
Semantic classifications and technical nouns are controlled terms; conformance depends on using them precisely.

\standard{This clause belongs to ISO/IEC 14882:2024; its precise constraints must be verified in the standard, including subclauses and any cross references.}
\textbf{Common trap.} Calling unspecified behavior undefined behavior, or vice versa.

\textbf{Self-test.} What demands documentation from an implementation?

\textbf{Answer rationale.} Implementation-defined behavior; unspecified behavior does not necessarily.

\practice{Make a one-page note with the relevant stable names, one original well-formed example, a negative or edge case, and a recorded toolchain assumption.}
\section{Clause 4: General principles}
\anchorref{4}
Implementation compliance, hosted/freestanding models and the abstract machine establish the interpretation framework.

\standard{This clause belongs to ISO/IEC 14882:2024; its precise constraints must be verified in the standard, including subclauses and any cross references.}
\textbf{Common trap.} Treating successful compilation as proof of defined behavior.

\textbf{Self-test.} What must happen for an ordinary diagnosable violation?

\textbf{Answer rationale.} A conforming implementation must issue at least one diagnostic.

\practice{Make a one-page note with the relevant stable names, one original well-formed example, a negative or edge case, and a recorded toolchain assumption.}
\section{Clause 5: Lexical conventions}
\anchorref{5}
Translation phases, tokenization, literals, identifiers and character sets govern source interpretation.

\standard{This clause belongs to ISO/IEC 14882:2024; its precise constraints must be verified in the standard, including subclauses and any cross references.}
\textbf{Common trap.} Assuming tokens are already final before preprocessing.

\textbf{Self-test.} Why might an identifier change validity in C++23?

\textbf{Answer rationale.} Unicode identifier constraints and normalization requirements.

\practice{Make a one-page note with the relevant stable names, one original well-formed example, a negative or edge case, and a recorded toolchain assumption.}
\section{Clause 6: Basics}
\anchorref{6}
Definitions, scopes, lookup, object lifetime, type properties, and execution relationships underlie all higher-level features.

\standard{This clause belongs to ISO/IEC 14882:2024; its precise constraints must be verified in the standard, including subclauses and any cross references.}
\textbf{Common trap.} Relying on a dangling object despite retaining its pointer value.

\textbf{Self-test.} What is the lifetime question for a view?

\textbf{Answer rationale.} Whether referenced objects remain alive and accessible at use time.

\practice{Make a one-page note with the relevant stable names, one original well-formed example, a negative or edge case, and a recorded toolchain assumption.}
\section{Clause 7: Expressions}
\anchorref{7}
Value categories, conversions, sequencing, operator behavior and constant evaluation affect meaning and overload selection.

\standard{This clause belongs to ISO/IEC 14882:2024; its precise constraints must be verified in the standard, including subclauses and any cross references.}
\textbf{Common trap.} Assuming left-to-right evaluation of all call arguments.

\textbf{Self-test.} Why can simplifying an expression alter overload resolution?

\textbf{Answer rationale.} It can change value category, type or conversion sequence.

\practice{Make a one-page note with the relevant stable names, one original well-formed example, a negative or edge case, and a recorded toolchain assumption.}
\section{Clause 8: Statements}
\anchorref{8}
Control flow includes selection, iteration, jumps and range-for semantics, including lifetime extensions.

\standard{This clause belongs to ISO/IEC 14882:2024; its precise constraints must be verified in the standard, including subclauses and any cross references.}
\textbf{Common trap.} Treating a loop initializer temporary exactly as in C++20.

\textbf{Self-test.} Which C++23 change affects temporary lifetime in range-for?

\textbf{Answer rationale.} Certain temporaries in the for-range-initializer live through the loop.

\practice{Make a one-page note with the relevant stable names, one original well-formed example, a negative or edge case, and a recorded toolchain assumption.}
\section{Clause 9: Declarations}
\anchorref{9}
Specifiers, declarators, initialization, function definitions, attributes, structured bindings and namespaces define entities.

\standard{This clause belongs to ISO/IEC 14882:2024; its precise constraints must be verified in the standard, including subclauses and any cross references.}
\textbf{Common trap.} Confusing direct/list initialization, or assuming default initialization initializes a scalar.

\textbf{Self-test.} Why can braced initialization change selected constructor?

\textbf{Answer rationale.} Initializer-list preference and narrowing restrictions affect selection.

\practice{Make a one-page note with the relevant stable names, one original well-formed example, a negative or edge case, and a recorded toolchain assumption.}
\section{Clause 10: Modules}
\anchorref{10}
Module units, imports, exports, fragments and reachability delimit compiled interfaces.

\standard{This clause belongs to ISO/IEC 14882:2024; its precise constraints must be verified in the standard, including subclauses and any cross references.}
\textbf{Common trap.} Assuming an exported name is automatically built or available in every toolchain.

\textbf{Self-test.} Which concern is language-defined and which is vendor-defined?

\textbf{Answer rationale.} Reachability semantics are standardized; BMI invocation is vendor-defined.

\practice{Make a one-page note with the relevant stable names, one original well-formed example, a negative or edge case, and a recorded toolchain assumption.}
\section{Clause 11: Classes}
\anchorref{11}
Special member functions, base subobjects, member access, initialization and comparisons define class semantics.

\standard{This clause belongs to ISO/IEC 14882:2024; its precise constraints must be verified in the standard, including subclauses and any cross references.}
\textbf{Common trap.} Breaking an invariant after move or during constructor failure.

\textbf{Self-test.} What is the rule-of-zero motivation?

\textbf{Answer rationale.} Delegate ownership to members whose special members are correct.

\practice{Make a one-page note with the relevant stable names, one original well-formed example, a negative or edge case, and a recorded toolchain assumption.}
\section{Clause 12: Overloading}
\anchorref{12}
Candidate formation, viability, best viable function and conversions determine call resolution.

\standard{This clause belongs to ISO/IEC 14882:2024; its precise constraints must be verified in the standard, including subclauses and any cross references.}
\textbf{Common trap.} Expecting an added overload always to be backward compatible.

\textbf{Self-test.} What can adding an overload do?

\textbf{Answer rationale.} It may create ambiguity or change the selected function.

\practice{Make a one-page note with the relevant stable names, one original well-formed example, a negative or edge case, and a recorded toolchain assumption.}
\section{Clause 13: Templates}
\anchorref{13}
Parameter substitution, constraints, lookup, instantiation and specialization define generic code.

\standard{This clause belongs to ISO/IEC 14882:2024; its precise constraints must be verified in the standard, including subclauses and any cross references.}
\textbf{Common trap.} Assuming syntactic concept satisfaction proves all semantic requirements.

\textbf{Self-test.} What is the distinction between satisfies and models?

\textbf{Answer rationale.} Modeling additionally requires the semantic conditions of the concept.

\practice{Make a one-page note with the relevant stable names, one original well-formed example, a negative or edge case, and a recorded toolchain assumption.}
\section{Clause 14: Exception handling}
\anchorref{14}
Throwing, unwinding, handlers, exception specifications and special exception-related functions define failures.

\standard{This clause belongs to ISO/IEC 14882:2024; its precise constraints must be verified in the standard, including subclauses and any cross references.}
\textbf{Common trap.} Calling noexcept a guarantee that no failure occurs.

\textbf{Self-test.} What happens when an exception escapes a noexcept(true) function?

\textbf{Answer rationale.} Termination is invoked.

\practice{Make a one-page note with the relevant stable names, one original well-formed example, a negative or edge case, and a recorded toolchain assumption.}
\section{Clause 15: Preprocessing directives}
\anchorref{15}
Conditional inclusion, macros, header imports and predefined names operate before core translation semantics.

\standard{This clause belongs to ISO/IEC 14882:2024; its precise constraints must be verified in the standard, including subclauses and any cross references.}
\textbf{Common trap.} Feature-detecting a library using only the language dialect flag.

\textbf{Self-test.} What C++23 directives simplify conditional compilation?

\textbf{Answer rationale.} \#elifdef and \#elifndef.

\practice{Make a one-page note with the relevant stable names, one original well-formed example, a negative or edge case, and a recorded toolchain assumption.}
\section{Clause 16: Library introduction}
\anchorref{16}
Header structure, requirements, freestanding subsets, library-wide contracts and implementation freedom establish library semantics.

\standard{This clause belongs to ISO/IEC 14882:2024; its precise constraints must be verified in the standard, including subclauses and any cross references.}
\textbf{Common trap.} Assuming a missing C++23 header means the compiler language parser is nonconforming.

\textbf{Self-test.} Why test compiler and library independently?

\textbf{Answer rationale.} They have distinct implementation and deployment capabilities.

\practice{Make a one-page note with the relevant stable names, one original well-formed example, a negative or edge case, and a recorded toolchain assumption.}
\section{Clause 17: Language support library}
\anchorref{17}
Runtime support, source locations, arithmetic properties, exceptions, coroutines and dynamic storage services support programs.

\standard{This clause belongs to ISO/IEC 14882:2024; its precise constraints must be verified in the standard, including subclauses and any cross references.}
\textbf{Common trap.} Treating source\_location strings as authenticated source identity.

\textbf{Self-test.} What does source\_location represent?

\textbf{Answer rationale.} Source location information determined by the compiler/context.

\practice{Make a one-page note with the relevant stable names, one original well-formed example, a negative or edge case, and a recorded toolchain assumption.}
\section{Clause 18: Concepts library}
\anchorref{18}
Concepts such as regular, same\_as, integral and invocable constrain generic algorithms.

\standard{This clause belongs to ISO/IEC 14882:2024; its precise constraints must be verified in the standard, including subclauses and any cross references.}
\textbf{Common trap.} Ignoring semantic axioms like equality preservation.

\textbf{Self-test.} What additional property distinguishes semantic modeling?

\textbf{Answer rationale.} Required semantic behavior beyond syntactic satisfaction.

\practice{Make a one-page note with the relevant stable names, one original well-formed example, a negative or edge case, and a recorded toolchain assumption.}
\section{Clause 19: Diagnostics library}
\anchorref{19}
Errors, assertions, system errors and stacktrace support diagnosis and failure reporting.

\standard{This clause belongs to ISO/IEC 14882:2024; its precise constraints must be verified in the standard, including subclauses and any cross references.}
\textbf{Common trap.} Using assertions as a production validation substitute.

\textbf{Self-test.} What distinguishes error\_code from throwing an exception?

\textbf{Answer rationale.} An explicit error object instead of control-flow unwinding.

\practice{Make a one-page note with the relevant stable names, one original well-formed example, a negative or edge case, and a recorded toolchain assumption.}
\section{Clause 20: Memory management library}
\anchorref{20}
Smart pointers, allocators and polymorphic memory resources manage ownership and allocation.

\standard{This clause belongs to ISO/IEC 14882:2024; its precise constraints must be verified in the standard, including subclauses and any cross references.}
\textbf{Common trap.} Assuming shared\_ptr prevents logical cycles or data races.

\textbf{Self-test.} When would weak\_ptr be important?

\textbf{Answer rationale.} Breaking non-owning observation out of shared-ownership cycles.

\practice{Make a one-page note with the relevant stable names, one original well-formed example, a negative or edge case, and a recorded toolchain assumption.}
\section{Clause 21: Metaprogramming library}
\anchorref{21}
Type traits, integer sequences and compile-time utilities support generic programming.

\standard{This clause belongs to ISO/IEC 14882:2024; its precise constraints must be verified in the standard, including subclauses and any cross references.}
\textbf{Common trap.} Assuming a false trait implies runtime behavior is always safe.

\textbf{Self-test.} When should a trait be consulted?

\textbf{Answer rationale.} When it provides a compile-time property required by an operation.

\practice{Make a one-page note with the relevant stable names, one original well-formed example, a negative or edge case, and a recorded toolchain assumption.}
\section{Clause 22: General utilities library}
\anchorref{22}
optional, variant, expected, function objects, formatting and numeric utilities support common abstractions.

\standard{This clause belongs to ISO/IEC 14882:2024; its precise constraints must be verified in the standard, including subclauses and any cross references.}
\textbf{Common trap.} Reading expected value when error state is active.

\textbf{Self-test.} When should expected be preferred to optional?

\textbf{Answer rationale.} When the failure alternative carries meaningful error information.

\practice{Make a one-page note with the relevant stable names, one original well-formed example, a negative or edge case, and a recorded toolchain assumption.}
\section{Clause 23: Strings library}
\anchorref{23}
Owned strings, string\_view and null-terminated utilities have different lifetime and format contracts.

\standard{This clause belongs to ISO/IEC 14882:2024; its precise constraints must be verified in the standard, including subclauses and any cross references.}
\textbf{Common trap.} Returning string\_view backed by a temporary.

\textbf{Self-test.} How can a string\_view be valid?

\textbf{Answer rationale.} Its referenced storage exists and covers the view range.

\practice{Make a one-page note with the relevant stable names, one original well-formed example, a negative or edge case, and a recorded toolchain assumption.}
\section{Clause 24: Containers library}
\anchorref{24}
Container requirements, invalidation, adaptors and multidimensional views govern sequence and associative storage.

\standard{This clause belongs to ISO/IEC 14882:2024; its precise constraints must be verified in the standard, including subclauses and any cross references.}
\textbf{Common trap.} Assuming vector iterators survive growth or erase.

\textbf{Self-test.} What must be reviewed after push\_back?

\textbf{Answer rationale.} Iterator/reference invalidation and reallocation conditions.

\practice{Make a one-page note with the relevant stable names, one original well-formed example, a negative or edge case, and a recorded toolchain assumption.}
\section{Clause 25: Iterators library}
\anchorref{25}
Iterator categories and concepts establish traversal and algorithm compatibility.

\standard{This clause belongs to ISO/IEC 14882:2024; its precise constraints must be verified in the standard, including subclauses and any cross references.}
\textbf{Common trap.} Reusing a single-pass input iterator as if it were forward.

\textbf{Self-test.} Why do iterator concepts matter?

\textbf{Answer rationale.} They define supported operations and multi-pass guarantees.

\practice{Make a one-page note with the relevant stable names, one original well-formed example, a negative or edge case, and a recorded toolchain assumption.}
\section{Clause 26: Ranges library}
\anchorref{26}
Range access, views, factories, adaptors and coroutine generators support compositional traversal.

\standard{This clause belongs to ISO/IEC 14882:2024; its precise constraints must be verified in the standard, including subclauses and any cross references.}
\textbf{Common trap.} Assuming a lazy view owns its backing sequence.

\textbf{Self-test.} What should be checked at every adapter boundary?

\textbf{Answer rationale.} Borrowing/lifetime and iterator semantics.

\practice{Make a one-page note with the relevant stable names, one original well-formed example, a negative or edge case, and a recorded toolchain assumption.}
\section{Clause 27: Algorithms library}
\anchorref{27}
Algorithm requirements, modifiers, sorting, numeric algorithms and parallel execution impose contracts and complexity.

\standard{This clause belongs to ISO/IEC 14882:2024; its precise constraints must be verified in the standard, including subclauses and any cross references.}
\textbf{Common trap.} Using an invalid comparator that lacks strict weak ordering.

\textbf{Self-test.} What sorting relation property is central?

\textbf{Answer rationale.} Strict weak ordering for appropriate comparator-based sorts.

\practice{Make a one-page note with the relevant stable names, one original well-formed example, a negative or edge case, and a recorded toolchain assumption.}
\section{Clause 28: Numerics library}
\anchorref{28}
Random facilities, numerical operations and floating point interfaces define mathematical computation semantics.

\standard{This clause belongs to ISO/IEC 14882:2024; its precise constraints must be verified in the standard, including subclauses and any cross references.}
\textbf{Common trap.} Using a pseudo-random engine as an entropy source for secrets.

\textbf{Self-test.} Why do uniform random engines not automatically give cryptographic secrecy?

\textbf{Answer rationale.} They do not specify a cryptographic security property.

\practice{Make a one-page note with the relevant stable names, one original well-formed example, a negative or edge case, and a recorded toolchain assumption.}
\section{Clause 29: Time library}
\anchorref{29}
Durations, time points, calendars and time zones distinguish clocks and civil time.

\standard{This clause belongs to ISO/IEC 14882:2024; its precise constraints must be verified in the standard, including subclauses and any cross references.}
\textbf{Common trap.} Using system\_clock to measure monotonic elapsed time.

\textbf{Self-test.} Which clock usually suits interval measurement?

\textbf{Answer rationale.} steady\_clock.

\practice{Make a one-page note with the relevant stable names, one original well-formed example, a negative or edge case, and a recorded toolchain assumption.}
\section{Clause 30: Localization library}
\anchorref{30}
Locales and facets influence culturally dependent formatting and classification.

\standard{This clause belongs to ISO/IEC 14882:2024; its precise constraints must be verified in the standard, including subclauses and any cross references.}
\textbf{Common trap.} Allowing local number formatting to change a machine-readable protocol.

\textbf{Self-test.} What locale should a wire format use?

\textbf{Answer rationale.} An explicitly defined protocol representation, often locale-independent.

\practice{Make a one-page note with the relevant stable names, one original well-formed example, a negative or edge case, and a recorded toolchain assumption.}
\section{Clause 31: Input/output library}
\anchorref{31}
Streams, filesystem, spans, formatting, printing and buffering define IO semantics.

\standard{This clause belongs to ISO/IEC 14882:2024; its precise constraints must be verified in the standard, including subclauses and any cross references.}
\textbf{Common trap.} Assuming filesystem path canonicalization is security authorization.

\textbf{Self-test.} What must an IO error path check?

\textbf{Answer rationale.} Stream state or error\_code/exception and resource invariants.

\practice{Make a one-page note with the relevant stable names, one original well-formed example, a negative or edge case, and a recorded toolchain assumption.}
\section{Clause 32: Regular expressions library}
\anchorref{32}
The modified ECMAScript regex grammar, iterators and match algorithms specify pattern processing.

\standard{This clause belongs to ISO/IEC 14882:2024; its precise constraints must be verified in the standard, including subclauses and any cross references.}
\textbf{Common trap.} Assuming a standards-compliant regex is safe against costly adversarial input.

\textbf{Self-test.} What is a regex resource test?

\textbf{Answer rationale.} A bounded-input adversarial runtime and memory evaluation.

\practice{Make a one-page note with the relevant stable names, one original well-formed example, a negative or edge case, and a recorded toolchain assumption.}
\section{Clause 33: Concurrency support library}
\anchorref{33}
Threads, atomics, synchronization, semaphores, futures and coordination types formalize shared execution.

\standard{This clause belongs to ISO/IEC 14882:2024; its precise constraints must be verified in the standard, including subclauses and any cross references.}
\textbf{Common trap.} Assuming relaxed atomic operations publish adjacent non-atomic data.

\textbf{Self-test.} What must be proven for shared non-atomic state?

\textbf{Answer rationale.} An adequate happens-before relationship or mutual exclusion.

\practice{Make a one-page note with the relevant stable names, one original well-formed example, a negative or edge case, and a recorded toolchain assumption.}
\section{Deep-dive exercise: a lifetime and container bug}
\begin{lstlisting}[language=C++]
#include <string>
#include <string_view>
std::string_view bad_label() {
    std::string s{"temporary"};
    return s; // dangling view after return
}
std::string good_label() {
    return std::string{"owned"};
}
\end{lstlisting}
Explain why returning a pointer-like view does not extend object lifetime. Link the failure to Clause 6 (object lifetime), Clause 23 (string view semantics), and Clause 20 where ownership choices are relevant. Repair the interface and write a regression test with suitable instrumentation.

\section{Deep-dive exercise: error model vs exception model}
\begin{lstlisting}[language=C++]
#include <expected>
#include <string_view>
enum class Error { invalid };
std::expected<int, Error> parse_one(std::string_view s) {
    if (s.size() != 1 || s[0] < '0' || s[0] > '9')
        return std::unexpected(Error::invalid);
    return s[0] - '0';
}
\end{lstlisting}
Specify the function precondition, its result alternatives, the possible allocation/exception model, and how callers propagate the error. Connect this exercise to Clauses 14, 19 and 22; discuss why expected does not prevent all exceptions.

\section{Deep-dive exercise: overload viability}
\begin{lstlisting}[language=C++]
void f(int);
void f(double);
// f(1) selects f(int); f(1.0) selects f(double).
// Add f(long), then test conversions for each call site.
// Test an intentionally ambiguous call separately.
\end{lstlisting}
Using Clauses 7, 9, 12 and 13, explain candidate sets, conversion sequence ranking and the effect of constraints. Avoid making claims about a rejected call without identifying whether the rule is diagnosable.

\section{Answer certification rubric}
\small
\begin{longtable}{@{}p{.18\linewidth}p{.40\linewidth}p{.34\linewidth}@{}}
\toprule
\textbf{Criterion} & \textbf{Points} & \textbf{What earns full credit}\\
\midrule
\endfirsthead
\toprule
\textbf{Criterion} & \textbf{Points} & \textbf{What earns full credit}\\
\midrule
\endhead
Vocabulary and status & 20 & Correct distinctions among undefined/unspecified/implementation-defined behavior\\
\addlinespace[2pt]
Clause trace & 20 & Accurate numbered and stable-name locators\\
\addlinespace[2pt]
Formal reasoning & 25 & Assumptions, semantics and conclusion are explicit\\
\addlinespace[2pt]
Executable evidence & 20 & Valid test with reproducible compiler command\\
\addlinespace[2pt]
Portability & 15 & Vendor/library limits and behavior categories documented\\
\addlinespace[2pt]
\bottomrule
\end{longtable}
\normalsize
\docend
\end{document}
